\chapter{INTRODUCTION}

\section{Background of research}\label{sec:intro-background}

Bike-sharing systems (BSSs) have been deployed in many cities worldwide as a flexible form of urban mobility and as part of environmentally sustainable transport systems \citep{alvarez-valdesOptimizingLevelService2016}. As of the end of 2024, there were 9,278,605 shared bicycles operating in 2,145 systems across over 1,700 cities in 92 countries \citep{bswmreport2024}. As a substitute for motorized travel, bike sharing helps relieve traffic congestion and reduce greenhouse gas and pollutant emissions \citep{luo2019comparativelife,kou2020quantifyinggreenhouse}, and it is well suited to short trips and to connections with public transit, including first- and last-mile travel \citep{nathModellingMethodsPlanning2019,demaioBikesharingHistoryImpacts2009,jin2025connectinglast}.

Based on how shared bicycles are locked and parked, BSSs are commonly classified into station-based systems, where bicycles are parked at designated docking stations, and free-floating systems, where bicycles are equipped with automatic locks and can be parked relatively freely. This thesis focuses on station-based systems. In a station-based BSS, a user rents a bicycle at one station and returns it at another. Consequently, a bike-sharing trip can fail in two ways: the user finds no bicycle to rent at the origin, or no vacant locker in which to return one at the destination \citep{alvarez-valdesOptimizingLevelService2016,ravivStaticRepositioningBikesharing2013}.

These failures arise because rental and return demand vary across stations and over time \citep{datnerSettingInventoryLevels2019}. The resulting spatio-temporal imbalance drives some stations toward a shortage of bicycles and others toward a shortage of vacant lockers, leaving arriving users unable to complete their trips \citep{brinkmannInventoryRoutingBike2016}. To reduce the service failures, operators redistribute bicycles among stations using a fleet of dedicated vehicles, an operation known as bike repositioning \citep{ravivStaticRepositioningBikesharing2013}. Generally, bike repositioning can be divided into two types: static and dynamic repositioning. Static repositioning is carried out during an off-peak period, typically overnight when the system is nearly idle and demand during the operation is negligible \citep{ravivStaticRepositioningBikesharing2013}. Dynamic repositioning is instead performed during the day to react to evolving demand \citep{caggianiModelingFrameworkDynamic2018}. The problem of determining the route of each vehicle together with the quantities of bicycles loaded or unloaded at the visited stations is known as the bike repositioning problem (BRP) \citep{ravivStaticRepositioningBikesharing2013,brinkmannInventoryRoutingBike2016}. In particular, this thesis studies the static BRP.

In addition to the usable bicycles handled in conventional repositioning, operators also need to deal with broken bicycles, which cannot be rented to satisfy rental demand. Meanwhile, a broken bicycle parked at a station occupies a locker that could otherwise accommodate a returned bicycle. Broken bicycles are a recurring concern across BSSs worldwide. In Barcelona, about 12\% of the bicycles available at stations were reported to have major defects that prevented normal use, and a further 55\% had minor defects such as broken bells and lights \citep{castrofernandezContributionBikesharingSustainable2011}. Similar patterns are observed elsewhere: an inspection of BIXI Montreal found roughly one in five bicycles in a state of disrepair \citep{demaioMeddinBikeSharingBlogMontreal2009}, while in New York's Citi Bike system about 2\% of the fleet requires major repairs every day \citep{kaspiBikesharingSystemsUser2017}.

To handle broken bicycles, operators commonly collect them and transport them to a depot for repair, thereby freeing occupied lockers at stations. Some broken bicycles may also be repaired on site as part of routine maintenance, restoring them to service without depot transport. For example, Vélib' in Paris employs maintenance agents who repair bicycles on the street rather than transporting them to repair shops \citep{tironiRepairingMaintainingPublic2012}, and Bicing in Barcelona deploys rapid-intervention motorbikes for minor repairs at stations \citep{clearchannelBicingBarcelona2010}. In practice, therefore, operators need to coordinate the repositioning of usable bicycles with the collection of broken bicycles and, where possible, their repair on site.

These operations jointly determine the resulting numbers of usable bicycles and vacant lockers at stations for the next day, and hence the system's ability to satisfy subsequent rental and return demand. In real systems, such decisions are made for large networks. In April 2026, for example, New York City's Citi Bike operated about 2,352 active docking stations and a fleet of more than 35,000 bicycles \citep{citibikeAprilMonthlyReport2026}. Coordinating the repositioning of usable bicycles with the collection and on-site repair of broken ones, at this scale and under demand that varies across stations and over time, poses a substantial operational challenge for system operators.

\section{Motivations and objectives}\label{sec:intro-motivations}

Quantifying how well a station will serve its users on the next day requires a measure of the dissatisfaction that arises when its bicycles or lockers run short. In this thesis, such evaluation is operationalized through two service-level measures. The user dissatisfaction function (UDF) \citep{ravivOptimalInventoryManagement2013} expresses the expected shortages of bicycles and lockers at a station as a function of its initial inventory, while the extended user dissatisfaction function (EUDF) \citep{kaspiBikesharingSystemsUser2017} further accounts for the initial inventory of broken bicycles. Based on these service-level measures, this thesis is motivated by three research gaps in static bike repositioning.

First, UDF-based static bike repositioning remains difficult to solve at large station-network scales. The problem jointly determines station-level bicycle movements and vehicle routes, and its search space grows rapidly with the numbers of stations and vehicles. Clustering-based decomposition can ease this burden by reducing the network size before detailed routing is performed \citep{forma3stepMathHeuristic2015}. Existing decompositions form clusters mainly from geographic and inventory considerations \citep{forma3stepMathHeuristic2015} or from service-level feasibility and approximate routing costs \citep{schuijbroekInventoryRebalancingVehicle2017}. The estimated UDF reduction attainable within a cluster and the per-vehicle routing-time feasibility have not yet been considered together at the clustering stage, so the resulting partition may be misaligned with both service improvement and route feasibility.

Second, broken bicycles create additional operational work that needs to be planned jointly with usable-bike repositioning. Some broken bicycles may be collected by trucks and transported to a depot, while others may be repaired on site by repairers. This workload is substantial. Citi Bike's publicly available monthly operating reports record the number of bicycle repairs completed each month, which, aggregated by calendar year, give the annual repair activity shown in Figure~\ref{fig:ch1-citibike-repair-activity}.

\begin{figure}[htbp]
    \centering
    \includegraphics[width=0.92\textwidth]{Chapter1/citibike_annual_bicycle_repair_activity_2017_2025.pdf}
    \caption[Annual bicycle repair activity from 2017--2025]{Annual bicycle repair activity from 2017--2025. Data obtained from Citi Bike Monthly reports.\protect\footnotemark}
    \label{fig:ch1-citibike-repair-activity}
\end{figure}
\footnotetext{\url{https://citibikenyc.com/system-data/operating-reports}}

The reported annual repair activity rose from 74,464 completed bicycle repairs in 2017 to a peak of 594,649 in 2022 and remained substantial at 265,310 in 2025. The workload is also visible in recent monthly operations: Citi Bike reported 11,637 completed depot repairs and 10,565 unique bicycles checked or repaired across depot and field operations in April 2026 alone \citep{citibikeAprilMonthlyReport2026}. Handling this workload interacts directly with usable-bike redistribution. Depot collection frees the lockers occupied by broken bicycles, but it also consumes truck capacity that would otherwise serve usable-bike redistribution \citep{wangStaticGreenRepositioning2018}. Because the associated truck routes and loads affect both operating effort and emissions, broken-bicycle handling also raises an environmental concern. Frequent usable-bike repositioning can substantially increase the fuel consumption of repositioning vehicles, which conflicts with the environmental motivation for operating BSSs as an eco-friendly transport mode \citep{luo2019comparativelife}. In London, for example, fossil-fueled vehicle usage for repositioning has been reported to outweigh the environmental benefit of replacing car trips by bike-sharing trips \citep{fishmanBikeSharesImpact2014}. Since repositioning plans with long routes or heavy vehicle loads may increase air-pollutant emissions \citep{shuiDynamicGreenBike2018}, the collection of broken bicycles can further intensify this concern: collected broken bicycles occupy vehicle capacity and have to be carried until the trucks return to the depot. On-site repair introduces an additional intervention, in which repairers restore broken bicycles in the field. The resulting planning problem must coordinate truck-based usable-bike repositioning, truck-based broken-bicycle collection, and repairer-based on-site repair. Previous BRP studies have incorporated broken-bicycle collection into repositioning. However, coordinating usable-bike repositioning, broken-bicycle collection, and repairer-based on-site repair within a single static plan has, to the best of the author's knowledge, not yet been studied.

Third, the integrated problem with broken bicycles and on-site repair requires a scalable decomposition method of its own. The static plan described above is difficult to solve at large station-network scales, so stations must be grouped before detailed truck and repairer routes are optimized. At this clustering stage, however, a candidate cluster can no longer be evaluated only by one-dimensional usable-bike inventory changes. Its value depends on the bivariate state of each station, on the combined effect of usable-bike transfers, broken-bicycle collection, and on-site repair, and on whether the truck and repairer workloads, which proceed in parallel, both fit within a common operating time budget. Existing clustering-based decompositions have not evaluated clusters through bivariate EUDF changes under parallel truck and repairer time demands. This creates a distinct large-scale clustering problem in which service improvement and operational feasibility need to be evaluated over coupled usable-bike and broken-bicycle states.

To fill the abovementioned research gaps, this thesis aims to:

\begin{enumerate}
    \item To develop a UDF-aware cluster-first, route-second framework for the multi-vehicle static bike repositioning problem, in which estimated UDF reduction and per-vehicle routing-time feasibility are incorporated into the clustering stage, and to evaluate its UDF-based service performance and computational scalability on large station networks.

    \item To formulate and solve an integrated static bike repositioning problem that coordinates truck-based usable-bike repositioning and broken-bicycle collection with repairer-based on-site repair, and to examine the operational conditions under which introducing on-site repairers is cost-effective.

    \item To develop an EUDF-aware large-scale decomposition framework for static bike repositioning with broken bicycles and on-site repair, in which candidate clusters are evaluated through bivariate EUDF changes and through parallel truck and repairer time components under a common time budget, and to examine how jointly accounting for broken-bicycle handling and repairer operations shapes the resulting clusters and large-scale solution quality.
\end{enumerate}

\section{Scope and outline of the thesis}\label{sec:intro-scope-outline}

The thesis focuses on station-based bike-sharing systems and static, off-peak planning. It considers multi-vehicle truck operations for usable-bike repositioning and broken-bicycle collection, together with independently routed repairers who perform on-site repairs. The models and solution methods developed in this thesis combine user-dissatisfaction-based service evaluation, mixed-integer linear programming, cluster-first, route-second decomposition, constructive heuristics, and hybrid genetic search.

The remainder of this thesis is organized as follows.

Chapter 2 reviews the literature on static bike repositioning, user-dissatisfac\-tion-based service-level evaluation, broken-bicycle collection, on-site repair operations, and decomposition methods for large-scale bike repositioning problems. It identifies the methodological and operational gaps addressed by the three research objectives.

Chapter 3 develops a UDF-aware cluster-first, route-second framework for the multi-vehicle static bike repositioning problem. It formulates the clustering stage to account for estimated UDF reduction and per-vehicle routing-time feasibility, introduces Activated Transfer Selection (ATS) to evaluate candidate clusters, and develops Hybrid Genetic Search for UDF-aware Clustering (HGS-UC) to solve the large-scale clustering problem.

Chapter 4 studies integrated usable-bike repositioning, broken-bicycle collection, and on-site repair. It presents a time-in\-dexed mixed-integer linear programming formulation that coordinates truck routes, repairer routes, station-level bike-handling quantities, and repair quantities. With truck routes and loads explicitly modeled, the objective also accounts for truck carbon emissions. It develops HGSADC-SBC, which combines Hybrid Genetic Search with Adaptive Diversity Control and a Sta\-tion-Budget-Constrained heuristic, as the solution method. The numerical analyses examine the operational conditions under which introducing on-site repairers is cost-effective.

Chapter 5 develops an EUDF-aware decomposition framework for large-scale static bike repositioning with broken-bicycle collection and on-site repair. It introduces Extended Activated Transfer Selection (EATS) to estimate cluster-level EUDF reduction while accounting for separate truck and repairer time requirements under a common operating time budget, and Hybrid Genetic Search for Extended-UDF Clustering (HGS-EUC) to solve the resulting large-scale station-partitioning problem. The cluster-level planning instances are then solved by the integrated solver developed in Chapter 4.

Chapter 6 summarizes the principal findings and contributions of the thesis, discusses its limitations, and outlines directions for future research.
